\chapter{Polyxan Hypergraphs: Definition and Basic Structure}
\label{ch:polyxan-basic}

Part~I established the geometric substrate of RSVP field theory on the semantic
manifold $(\mathcal{M},g)$.  
Part~II develops the algebraic and combinatorial substrate: the
\emph{Polyxan hypergraphs} that represent content atoms, typed links, media
structures, and their associated sheaf-data.

This chapter provides the rigorous foundation for all subsequent developments:
typed vertices, typed hyperedges, curvature invariants, morphisms, and the
construction of the category $\Polyx$, ultimately shown to be a presheaf topos.
It closes with the definition of a \emph{media quine}, the central fixed point
object of the extraction--hoisting--reduction pipeline.

% ===============================================================
\section{Polyxan Atoms and Typed Vertex Structure}
% ===============================================================

A Polyxan hypergraph is built from \emph{Polyxan atoms}, the fundamental
syntactic-semantic units of the theory.

\begin{definition}[Polyxan atom]
A \emph{Polyxan atom} is a tuple
\[
a = (t(a),\;\sigma(a),\;\tau(a),\;w(a)),
\]
where:
\begin{enumerate}
  \item $t(a)$ is a \emph{type label} drawn from a small category $\mathcal{T}$
        of semantic types (e.g.\ \textsc{concept}, \textsc{token},
        \textsc{relation}, \textsc{operation});
  \item $\sigma(a)$ is a \emph{content payload} (a finite structured object
        drawn from an ambient semantic set $S$);
  \item $\tau(a)$ is a \emph{semantic tag}, a finite tuple of invariants used to
        track modality, provenance, or context;
  \item $w(a)\in\mathbb{R}_{\ge 0}$ is a \emph{weight}, interpreted as
        semantic density or importance.
\end{enumerate}
We denote the set of all Polyxan atoms of a hypergraph $G$ as $V(G)$.
\end{definition}

Types are required to form at least a \emph{finitary} category---finite
objects, finite hom-sets---because the hypergraph rewriting systems of
Chapters~\ref{ch:polyxan-rewriting}--\ref{ch:polyxan-functorial} rely on
combinatorial finiteness conditions.

Tags provide a weak form of modality, covering distinctions such as textual,
visual, auditory, procedural, or hybrid content, but no modality is privileged.
Weights will feed into curvature measures in later sections.

% ===============================================================
\section{Typed Hyperedges and Source--Target Signatures}
% ===============================================================

Vertices alone do not encode semantic structure; the essential structure lies in
the \emph{typed hyperedges}.

\begin{definition}[Typed hyperedge]
A \emph{typed hyperedge} $e$ of arity $(m,n)$ in a Polyxan hypergraph $G$ is a
tuple
\[
e = (s_1,\ldots,s_m;\; t_1,\ldots,t_n;\; \lambda(e)),
\]
where:
\begin{enumerate}
  \item each $s_i,t_j \in V(G)$ are Polyxan atoms (sources and targets);
  \item $(s_1,\ldots,s_m)$ is the \emph{source signature};
  \item $(t_1,\ldots,t_n)$ is the \emph{target signature};
  \item $\lambda(e)$ is a \emph{hyperedge type} belonging to a finite set
        $\Lambda$ of link types with prescribed typing rules.
\end{enumerate}
We denote the set of all typed hyperedges of $G$ as $E(G)$.
\end{definition}

A typed hyperedge $e$ must satisfy \emph{type consistency}:
\[
t(s_i) \xrightarrow{\lambda(e)} t(t_j)
\quad
\text{in the type category } \mathcal{T}.
\]
Thus every hyperedge lifts canonically to a morphism in $\mathcal{T}$.

\begin{remark}[Multimodality and directionality]
Hyperedges may mix modalities (via tags) and are explicitly directed.
The $(m,n)$ arity need not be $(1,1)$; higher-arity interactions are essential.
\end{remark}

The incidence structure of a Polyxan hypergraph is thus much richer than that of
a simple graph or category: it includes multi-input, multi-output transformations
that correspond, in media contexts, to rewrites, operations, conversions,
translations, and multi-modal aggregations.

% ===============================================================
\section{Definition of a Polyxan Hypergraph}
% ===============================================================

\begin{definition}[Polyxan hypergraph]
A \emph{Polyxan hypergraph} $G$ consists of:
\begin{enumerate}
  \item a finite set $V(G)$ of Polyxan atoms;
  \item a finite set $E(G)$ of typed hyperedges;
  \item a sheaf of semantic data $\mathscr{F}_G$ over the incidence complex
        $\mathrm{Inc}(G)$ associating to each cell the admissible content
        restrictions and modality selectors;
  \item a curvature assignment $\kappa_G$ defined on vertices and hyperedges;
  \item gluing maps ensuring consistency on overlaps of the incidence complex.
\end{enumerate}
\end{definition}

\subsection{Incidence complex}

The incidence structure gives rise to a simplicial complex (or, more exactly, a
fat nerve) whose $k$-simplices correspond to chains of hyperedges.  The sheaf
$\mathscr{F}_G$ attaches admissible content constraints to simplices, enabling
context-sensitive interpretations, refinements, and rewrites.

\subsection{Curvature assignment}

The curvature $\kappa_G$ is built from three components:

\begin{enumerate}
  \item \emph{Hyperedge deficit:}  
        $\delta(e) = \mathrm{deg}(e) - \mathrm{rank}(e)$ measures how many
        combinatorial degrees of freedom a hyperedge consumes.
  \item \emph{Sectional curvature:}  
        defined on two-planes of the incidence complex using a synthetic metric.
  \item \emph{Sheaf curvature:}  
        measured by the failure of sheaf sections to glue coherently across
        triple overlaps.
\end{enumerate}

These curvature measures will serve as input for the Polyxan energy functional
and for the spectral geometry of hypergraphs (Chapter~\ref{ch:hypergraph-energy}).

% ===============================================================
\section{Morphisms and the Category \texorpdfstring{$\Polyx$}{Polyx}}
\label{sec:polyx-category}
% ===============================================================

Polyxan hypergraphs naturally organize into a category.

\begin{definition}[Morphisms]
A \emph{Polyxan morphism} $f : G \to H$ consists of:
\begin{enumerate}
  \item a type-preserving map $f_V : V(G) \to V(H)$ satisfying
        $t(f_V(a)) = t(a)$;
  \item an induced map $f_E : E(G) \to E(H)$ preserving source--target
        signatures and hyperedge types;
  \item a morphism of sheaves $\mathscr{F}_G \to f^\ast\mathscr{F}_H$ on the
        incidence complexes;
  \item preservation of curvature: $\kappa_H(f(x)) = \kappa_G(x)$ for all
        vertices and hyperedges $x$.
\end{enumerate}
Composition is defined componentwise.
\end{definition}

\begin{definition}[The category Polyx]
The category $\Polyx$ has:
\[
\mathrm{Obj}(\Polyx) = \{ \text{Polyxan hypergraphs} \},
\qquad
\mathrm{Mor}(\Polyx) = \{ \text{Polyxan morphisms} \}.
\]
\end{definition}

% ===============================================================
\section{Polyx as a Presheaf Topos}
% ===============================================================

We now demonstrate that $\Polyx$ is equivalent to a presheaf topos, giving it
an internal intuitionistic logic and enabling the powerful categorical
machinery of sheaf semantics.

\subsection{Site of semantic contexts}

Let $\mathcal{C}$ be the category whose objects are finite semantic contexts
(i.e.\ small finite type-labelled sets) and whose morphisms are context
refinements.  We equip $\mathcal{C}$ with the canonical Grothendieck topology
generated by covering families of refinements.

\begin{theorem}[Topos equivalence]
\label{thm:polyx-topos}
There is an equivalence of categories
\[
\Polyx \;\simeq\; \widehat{\mathcal{C}}
:= \mathrm{Funct}(\mathcal{C}^{op},\mathbf{Set}),
\]
showing that $\Polyx$ is a presheaf topos.
\end{theorem}

\begin{proof}[Sketch]
Every Polyxan hypergraph $G$ can be encoded as a contravariant functor
$\mathcal{C}^{op}\to\mathbf{Set}$ by sending a semantic context $U$ to the set
of all type-consistent embeddings of $U$ into $G$ respecting sheaf-data and
curvature.  Morphisms in $\Polyx$ induce natural transformations.  Full
faithfulness and essential surjectivity follow from the fact that contexts form
a basis for the semantic topology on hypergraphs, and gluing conditions for
sheaves guarantee the sheaf-compatible structure needed for the presheaf
form.
\end{proof}

Consequences include:

\begin{enumerate}
  \item $\Polyx$ has all limits and colimits;
  \item $\Polyx$ supports an internal intuitionistic logic;
  \item hypergraph rewriting can be interpreted as internal computation;
  \item the slice categories $\Polyx/G$ are themselves toposes.
\end{enumerate}

These results give Polyxan hypergraphs the same categorical robustness that
Part~I established for RSVP fields via their geometric foundations.

% ===============================================================
\section{Curvature Invariants of Polyxan Hypergraphs}
% ===============================================================

We now formalize the curvature measures introduced earlier.

\subsection{Hyperedge deficit}

For an edge $e$ with arity $(m,n)$ and rank 
$r(e)=\dim(\mathrm{Span}\{t(s_i)\to t(t_j)\})$,
define
\[
\delta(e) := (m+n) - r(e).
\]
Large deficit corresponds to combinatorial rigidity; small deficit corresponds
to combinatorial looseness.

\subsection{Sectional curvature}

Let $x\in G$ be a vertex or hyperedge.  
The local incidence complex gives a synthetic metric; the sectional curvature
$K(\Pi)$ of a 2-plane $\Pi$ at $x$ is computed analogously to CAT$(0)$
polyhedral geometry.

\subsection{Sheaf curvature}

Given three overlapping cells $U,V,W$ in $\mathrm{Inc}(G)$, define the sheaf curvature
\[
\Omega(U,V,W)
:= \dim \ker\!\big(
    \mathscr{F}_G(U)\oplus \mathscr{F}_G(V)\oplus \mathscr{F}_G(W)
    \to
    \mathscr{F}_G(U\cap V\cap W)
  \big).
\]
Nonzero $\Omega$ measures obstruction to gluing.

\subsection{Euler characteristic of the nerve}

The nerve $\mathcal{N}(G)$ of $G$ is defined by higher-order incidences.
Its Euler characteristic
\[
\chi(G) := \sum_{k\ge 0} (-1)^k |\mathcal{N}(G)_k|
\]
serves as a global combinatorial invariant, analogous to the topological Euler
characteristic of a manifold.

% ===============================================================
\section{Media Quines}
% ===============================================================

Polyxan hypergraphs evolve under extraction, hoisting, and reduction (developed
in detail in Chapters~\ref{ch:extract-hoist}--\ref{ch:polyquine-enum}).
A \emph{media quine} is a fixed point of this dynamic.

\begin{definition}[Media quine]
A Polyxan hypergraph $G$ is a \emph{media quine} if:
\begin{enumerate}
  \item extraction preserves $G$ up to isomorphism;
  \item hoisting preserves $G$ up to isomorphism;
  \item reduction preserves $G$ up to isomorphism;
  \item all curvature invariants $(\delta, K, \Omega, \chi)$ are preserved;
  \item the sheaf $\mathscr{F}_G$ is stable under all three transformations.
\end{enumerate}
\end{definition}

Media quines represent perfect self-sustaining semantic structures:  
closed under interpretation, compression, and re-expansion.  
They are the Polyxan analog of the semantic vacua of Part~I.

% ===============================================================
\section{Summary}
% ===============================================================

This chapter introduced Polyxan hypergraphs, their atoms, typed hyperedges,
sheaf data, curvature invariants, and their organization into the presheaf
topos $\Polyx$.  
It established the combinatorial and categorical foundation upon which the
remaining chapters of Part~II---rewriting systems, spectral semantics, and
functorial Polycompilation---will be built.

